%%%%%
%
%% Dans ce fichier, les �ventuels accents sont cod�s
%% en ISO-8859-15 (aka latin9)
%% La ligne suivante donne trois exemples 
%% e accent aigu : �, e accent grave : �, e-dans-l'o : �
\section{Calcul du cospectre de chaleur selon Moncrieff 1997}

\itemstyle{Op\'erateur~:} \operfmt{Cospectre\_Chaleur\_Moncrieff97}
\index{Cospectre\_Chaleur\_Moncrieff97}

%%%%%OBJET%%%%
\itemstyle{Objet~:}
Calcul du cospectre th�orique selon Moncrieff 1997 correspondant aux flux de chaleur sensible et latent.

%%%%%ARGUMENTS%%%%
\bigskip
%% Dans ce fichier, les �ventuels accents sont cod�s
%% en ISO-8859-15 (aka latin9)
%% La ligne suivante donne trois exemples 
%% e accent aigu : �, e accent grave : �, e-dans-l'o : �
%%%%%ENTREES%%%%
\begin{tabarg}
\multicolumn{3}{l}{\titretab {Entr\'ees}}\\
\hline
$N_{freq}$ & Coeff., \texttt{float} & Nombre de points du cospectre \\
$z$ & Coeff., \texttt{float} & Hauteur de mesure (m)\\
$d$ & Coeff., \texttt{float} & Hauteur de d�placement (m)\\
$L$ & Vecteur, \texttt{float} & Longueur de Monin-Obhukof (m)\\
$U_{moy}$ & Vecteur, \texttt{float} & Force du vent moyen (m/s)\\
\end{tabarg}

\medskip
%%%%%SORTIES%%%%
\begin{tabarg}
\multicolumn{3}{l}{\titretab {Sorties}}\\
\hline
$Cosp$ & Vecteur, \texttt{float[$N_{freq}$]} & Cospectre \\
\end{tabarg}


%%%%%FORMULE%%%%
\itemstyle{Formule}

%% Formule de l'operateur : Cospectre_Chaleur_Moncrieff97
%% Dans ce fichier, les �ventuels accents sont cod�s
%% en ISO-8859-15 (aka latin9)
%% La ligne suivante donne trois exemples 
%% e accent aigu : �, e accent grave : �, e-dans-l'o : �


\medskip

Pour j variant de 0 � $N_{freq} -1$, on calcule le tableau des fr�quences $f_j$  : 

$$
f_j = \exp(\ln(f_0)+ \frac{\ln(10.)-\ln(f_0)}{N_{freq}} * (j+1)
$$

avec $f_0 = 1/7200.$ \\


Pour chaque pas de temps $i$ on calcule :


\begin{eqnarray*}
k_{i,j} & = & f_j*\left\vert\frac{z - d}{U_{moy i}} \right\vert \\
\zeta_i & = & (z - d)/L_i
\end{eqnarray*}
     
Si $\zeta_i > 0. $ (conditions stables ) alors : 
\begin{eqnarray*}
 A_i & = & 0.284 * (1. + 6.4 * \zeta_i )^{0.75} \\
 B_i & = & 2.34 * A_i^{-1.1} \\
 Cosp_{i,j} & = & \frac{ k_ {i,j} }{ f_j * (A_i + B_i * f_j^{2.1}) }
\end{eqnarray*}
           
sinon (conditions instables et neutres)  : 


\hspace{3cm} si $k_{i,j} < 0.54$
\begin{eqnarray*}
Cosp_{i,j} & = & \frac{ 12.92 * k_{i,j} }{ f_j * (1. + 26.7 * k_{i,j})^{1.375} }
\end{eqnarray*}

\hspace{3cm} sinon :
\begin{eqnarray*}            
Cosp_{i,j} & = &  \frac{4.378 * k_{i,j} }{ f_j * (1. + 3.8 * k_{i,j})^{2.4} }
\end{eqnarray*}






\itemstyle{Remarques~:}
Si on ne connait pas la hauteur de d�placement $d$, on peut utiliser par d�faut $d=0.67*Hauteur\ de\ la\ vegetation$ (source EDDYPRO).

%%%%%AUTEUR%%%%
\itemstyle{Auteur~:}
CNRM/GMEI/TRAMM

%%%%%REFERENCE%%%%
\itemstyle{R\'ef\'erences~:}
\cite{MONCRIEFF-1997}
